\section{Introduction}

Text-to-speech (TTS) synthesis has undergone a fundamental shift from rule-based concatenative systems toward neural architectures
capable of producing perceptually indistinguishable speech. Within this trajectory, zero-shot voice cloning occupies a distinct
technical position: given only a brief reference recording---typically 20 to 30 seconds of natural speech---a model must condition
its generative process to reproduce the speaker's identity across arbitrary target text, without any speaker-specific fine-tuning.
The challenge is not simply one of signal reconstruction. It requires a generative model to disentangle speaker-characteristic
acoustic features (timbre, prosodic rhythm, vocal tract resonance) from linguistic content, encode them into a stable
representational space, and faithfully project them onto novel utterances the reference speaker never produced.

The emergence of large-scale neural codec language models \cite{wang2023neural, borsos2023audiolm}, flow-matching decoders
\cite{le2024voicebox}, and diffusion-based vocoders has produced a generation of zero-shot systems achieving mean opinion scores
competitive with ground-truth recordings in controlled English-language evaluations. Yet the field's evaluation practices remain
inconsistently scoped. Most published benchmarks operate on clean studio-quality reference audio, target a single high-resource
language, and report MOS under conditions that may not reflect deployment variability. The problem deepens in multilingual contexts.
Hindi, spoken by over 600 million people, presents phonological properties---retroflex consonants, vowel nasalization, and
Devanagari-to-phoneme mapping ambiguities---that stress both acoustic modeling and ASR-based intelligibility measurement in ways
English benchmarks do not reveal. A model that achieves low word error rate in English may degrade substantially when tasked with
synthesizing Hindi text, particularly if the architecture has not been trained on phonologically diverse corpora.

Speaker identity preservation compounds this difficulty. Reference audio clips of 20--30 seconds encode highly compressed
speaker-characteristic information; the model's speaker encoder must extract a stable embedding from this limited evidence and
maintain it across sentences of varying length and phonetic content. Cosine similarity computed over high-dimensional embedding
vectors (derived from independent speaker verification models such as WavLM or ECAPA-TDNN \cite{chen2022wavlm, desplanques2020ecapa})
provides a reproducible proxy for perceptual identity similarity, though it does not capture all dimensions of speaker recognition.
Acoustic naturalness---the degree to which synthesized speech is free of artifacts, unnatural prosody, or spectral discontinuities---
is equally difficult to operationalize at scale; neural MOS predictors such as UTMOS \cite{saeki2022utmos} enable non-intrusive
quality estimation without costly human listening panels, though their reliability across languages remains an active research question.

In this work, we conduct a systematic, multi-lingual empirical evaluation of three contemporary zero-shot voice cloning
architectures: \textbf{Orpheus-3B} \cite{orpheus2024}, a large language model-based TTS system adapted for zero-shot speaker
conditioning; \textbf{VoxtLM/VoxCPM2} \cite{voxcpm2024}, a codec-based generative model trained on diverse multilingual speech; and
\textbf{VibeVoice}, a diffusion-based synthesis pipeline with a contrastive speaker encoder. Each system receives an identical
20-to-30-second reference audio clip and a target text sequence; the synthesized output is evaluated across three complementary
dimensions: speaker similarity (cosine distance in speaker embedding space), intelligibility (Word Error Rate for English and
Character Error Rate for Hindi, both transcribed via Whisper ASR \cite{radford2023whisper}), and acoustic naturalness (UTMOS score).
The use of CER rather than WER for Hindi is deliberate---Whisper's Hindi transcription output exhibits transliteration
inconsistencies at the word boundary level, and character-level comparison isolates phoneme-fidelity from tokenization artifacts.

The significance of this comparison extends beyond benchmarking. Each architecture embodies different inductive biases regarding the
speaker conditioning mechanism, decoder design, and training regime. Orpheus-3B conditions generation autoregressively through
speaker-prefixed token sequences; VoxCPM2 operates through a codec language model with cross-attention to a speaker prototype;
VibeVoice uses a diffusion process guided by a contrastive speaker embedding. Identifying which of these mechanisms produces more
stable identity preservation under a low-data reference condition, and whether that stability transfers across linguistic
scripts, informs both model selection for deployment and the architectural choices most relevant for future multilingual TTS research.

The remainder of this paper is organized as follows: Section~\ref{sec:related} reviews related literature on neural TTS, zero-shot
speaker adaptation, and multilingual speech synthesis; Section~\ref{sec:methodology} details the experimental methodology,
including dataset construction, preprocessing pipeline, and evaluation metric implementation; Section~\ref{sec:results} presents
quantitative results across all three evaluation dimensions; Section~\ref{sec:discussion} analyzes performance trade-offs and
architectural correlates of observed differences; Section~\ref{sec:conclusion} summarizes findings and outlines directions for
future work.
